\subsection{微分方程的概念}

设 I 是包含于 R 中且不退化为一点的区间，E 是 R 上的一个拓扑向量空间，A 与 B 是 E 的两个开子集。设 \((\mathbf{x}, \mathbf{y}, t) \mapsto \mathbf{g}(\mathbf{x}, \mathbf{y}, t)\) 是 \(A \times B \times I\) 到 E 中的连续映射；对于 I 到 A 中的每个可导映射 u，若其导数取值于 B，我们使它对应 \(t \mapsto \mathbf{g}(\mathbf{u}(t), \mathbf{u}'(t), t)\) 这个 I 到 E 中的映射，并记之为 \(\tilde{\mathbf{g}}(\mathbf{u})\)；于是 \(\tilde{g}\) 定义在由 I 到 B 中、且导数取值于 B 的可导函数所成的集合 \(\mathcal{D}(\mathrm{A}, \mathrm{B})\) 上。我们称方程 \(\tilde{\mathbf{g}}(\mathbf{u}) = 0\) 为 u 的微分方程（关于实变量 t）；这个方程的解也称为微分方程（在区间 I 上的）积分；它是 I 到 A 中的一个可导映射，其导数取值于 B，且使得 \(\mathbf{g}(\mathbf{u}(t), \mathbf{u}'(t), t) = 0\) 对每个 \(t \in I\) 成立。滥用语词地说，我们把微分方程 \(\tilde{\mathbf{g}}(\mathbf{u}) = 0\) 写成如下形式

\[
\mathbf {g} (\mathbf {x}, \mathbf {x} ^ {\prime}, t) = 0,
\]

这里约定 x 属于集合 D(A, B)。

例如，当 I = E = R 时，关系式

\[
x ^ {\prime} = 2 t, \qquad t x ^ {\prime} - 2 x = 0, \qquad x ^ {\prime 2} - 4 x = 0, \qquad x - t ^ {2} = 0
\]

都是微分方程，这四个方程都以函数 \(x(t) = t^{2}\) 为解。

本章原则上只考虑 E 是 R 上的完备赋范空间、且微分方程具有下列特殊形式的情形

\[
\mathbf {x} ^ {\prime} = \mathbf {f} (t, \mathbf {x})\tag{1}
\]

（“关于 \(\mathbf{x}'\) 的显式方程”），其中 \(\mathbf{f}\) 表示定义在 \(\mathrm{I} \times \mathrm{H}\) 上、取值于 E 中的函数，而 H 是 E 的一个非空开子集。此外，我们还要把这类方程（在区间 I 上的）解（或积分）的概念稍加扩充：我们说，定义在 I 上、连续且取值于 H 中的函数 \(\mathbf{u}\)，若存在 I 的一个可数子集 A，使得在 A 关于 I 的补集的每一点 \(t\) 处函数 \(\mathbf{u}\) 都有导数 \(\mathbf{u}'(t)\)

使得 \(\mathbf{u}^{\prime}(t)=\mathbf{f}(t,\mathbf{u}(t))\)；若 u 可导且在每一点 \(t\in I\) 处满足上述关系，我们就称它为方程 (1) 在 I 上的严格解。

对于形如 \(\mathbf{x}^{\prime}=\mathbf{f}(t)\) 的微分方程这一特殊情形，其中 f 是 I 到 E 中的映射，上述意义下的解就是函数 f 的原函数（II，第 51 页），而严格解就是严格原函数。

当 E 是完备赋范空间 \(E_{i}\)（\(1 \leqslant i \leqslant n\)）的乘积时，可以写 \(\mathbf{x} = (\mathbf{x}_{i})_{1 \leqslant i \leqslant n}\) 与 \(\mathbf{f} = (\mathbf{f}_{i})_{1 \leqslant i \leqslant n}\)，其中 \(x_{i}\) 是 I 到 \(E_{i}\) 中的映射，\(f_{i}\) 是 \(I \times H\) 到 \(E_{i}\) 中的映射；于是方程 (1) 等价于微分方程组

\[
\mathbf {x} _ {i} ^ {\prime} = \mathbf {f} _ {i} (t, \mathbf {x} _ {1}, \mathbf {x} _ {2}, \dots , \mathbf {x} _ {n}) \quad (1 \leqslant i \leqslant n).\tag{2}
\]

最重要的情形是诸 \(E_{i}\) 都等于 R 或 C 的情形；这时称 (2) 为数量微分方程组。

可以把对形如

\[
\mathbf {x} ^ {(n)} = \mathbf {f} \big (t, \mathbf {x}, \mathbf {x} ^ {\prime}, \dots , \mathbf {x} ^ {(n - 1)} \big)\tag{3}
\]

的关系式的研究化归为对方程组 (2) 的研究，其中 \(\mathbf{x}\) 是 I 上 \(n\) 次可导的向量函数；因为若令 \(\mathbf{x}_1 = \mathbf{x}\)，并对 \(\mathbf{x}_p = \mathbf{x}^{(p - 1)}\) 令 \(2 \leqslant p \leqslant n\)，则关系式 (3) 等价于方程组

\[
\left\{ \begin{array}{l} \mathbf {x} _ {i} ^ {\prime} = \mathbf {x} _ {i + 1} \\ \mathbf {x} _ {n} ^ {\prime} = \mathbf {f} (t, \mathbf {x} _ {1}, \mathbf {x} _ {2}, \ldots , \mathbf {x} _ {n}). \end{array} \right. \quad (1 \leqslant i \leqslant n - 1)\tag{4}
\]

形如 (3) 的关系式称为 n 阶微分方程（已对 \(\mathbf{x}^{(n)}\) 明确解出）；相反，形如 (1) 的方程称为一阶微分方程。

类似地，可以把任何形如

\[
\mathrm{D} ^ {n _ {i}} \mathbf {x} _ {i} = \mathbf {f} _ {i} (t, \mathbf {x} _ {1}, \mathrm{D} \mathbf {x} _ {1}, \dots , \mathrm{D} ^ {n _ {1} - 1} \mathbf {x} _ {1}, \dots , \mathbf {x} _ {p}, \mathrm{D} \mathbf {x} _ {p}, \dots , \mathrm{D} ^ {n _ {p} - 1} \mathbf {x} _ {p})\tag{5}
\]

\((1 \leqslant i \leqslant p)\) 的“微分方程组”化归为形如 (2) 的方程组，其中 \(\mathbf{x}_i\) 是 I 上的函数，在 I 上 \(n_i\) 次可导（对 \(1 \leqslant i \leqslant p\)）。

